\section{Reflection}

\subsection{Working Inside Hytale's Development Pipeline}

Because I had previous Minecraft modding experience and a computer science background,
I expected \textit{Trial of the Goddess} to be challenging but manageable.
In practice,
Hytale's official modding pipeline was more restrictive than expected,
which surprised me because the game presents itself as a Minecraft alternative
built around modding,
creation,
and official creator tools.

One reason was that Hytale is still in active development.
During research,
I often found references to API calls,
classes,
or systems that had already been deprecated,
renamed,
removed,
or changed.
Since the game is still new,
there is also no large community knowledge base comparable to Minecraft modding.
Many problems could not simply be looked up;
I had to inspect asset files,
test systems myself,
and treat the existing game data as unofficial documentation.

The main difficulty was not simply writing Java code,
but understanding where Hytale allowed intervention.
Expected tasks such as commands,
chat messages,
custom UI,
item handling,
sound events,
and health changes worked reasonably well.
However,
the server-side plugin model created limits:
I could not freely change the client interface,
compass,
rendering behavior,
or every kind of player feedback.
When a needed system was not exposed,
I had to approximate it with existing tools
or redesign the feature.

The Goddess statue was the clearest example.
I first expected the interaction to work through Java event listeners.
Instead,
the statue only became usable after its block JSON pointed to my custom
\texttt{GoddessTrial\_Statue} interaction.
Asset files were therefore not only decorative resources,
but part of the logic that decided what the code could react to.
Much of development became searching existing assets with PowerShell commands
such as \texttt{Get-ChildItem} and \texttt{Get-Content}
to find examples of swords,
flowers,
furniture,
sound events,
interaction hints,
and block properties.

Several design choices came from these limits.
When something did not work as intended,
I first tried to work around it within Hytale's existing systems.
When a workaround became too unstable,
I changed the design instead of forcing the original plan.

The most disappointing example was the Sacred Flower.
I originally wanted it to spawn dynamically in random terrain
and developed a recursive positioner that would choose a nearby position
and iteratively move the flower until the placement checks were valid.
On paper,
this was the more interesting and mathematically satisfying solution.
In practice,
the terrain was too unpredictable:
the flower could appear in the air,
underground,
inside grass,
too far away,
or at invalid positions.
The system became unreliable,
caused crashes,
and took too much time to stabilize before the deadline.

I replaced it with manually defined positions in a fixed world.
This was less satisfying technically,
but it made the trial stable and playable.
From the player's perspective,
the flower still appears as the objective,
the player still has to find it,
and the ritual still works.
The same pattern appeared elsewhere:
the compass marker became the Sacred Flower's purple glow,
and recurring reinforcements were disabled because spawning NPCs from inside
a ticking system conflicted with Hytale's entity store processing.

These compromises were not simply failures.
They show that the development pipeline did not only enable the mod;
it also shaped what the mod could realistically become.
In the end,
I learned to read the game itself as its own documentation:
to follow the logic hidden in its assets,
to force a path through systems that were not always meant for me,
and to compromise without letting the central idea disappear.